\chapter{পর্যাবৃত্ত গতি}

\section{পর্যাবৃত্তি}
কোন কিছুর যদি নির্দিষ্ট স্থান বা কাল পর পর পুনরাবৃত্তি ঘটে তাকে পর্যাবৃত্তি (Periodicity) বলে। এমন ঘটনা, রাশি বা অপেক্ষককে পর্যাবৃত্তিক ঘটনা, রাশি বা অপেক্ষক বলে। পর্যাবৃত্তি দুই প্রকার যথা: স্থানিক পর্যাবৃত্তি ও কালিক পর্যাবৃত্তি।
পর্যাবৃত্তি নির্দিষ্ট দূরত্ব পর পর ঘটলে তাকে স্থানিক পর্যাবৃত্তি ও নির্দিষ্ট সময় পর পর ঘটলে তাকে কালিক পর্যাবৃত্তি বলে।

কোন বস্তু কণা এর গতিপথে কোন নির্দিষ্ট বিন্দুকে নির্দিষ্ট সময় পরপর একই দিক থেকে অতিক্রম করলে তাকে পর্যাবৃত্ত গতি (Periodic Motion) বলে। পর্যাবৃত্ত গতিসম্পন্ন কণা পর্যায়কালের অর্ধেক সময় সাম্যবস্থানের একদিকে এবং বাকি অর্ধেক সময় এক‌ই গতিপথে তার বিপরীত দিকে চললে তাকে স্পন্দন গতি/দোলন গতি (Harmonic Motion) বলে।

সরল দোলন গতি (Simple Harmonic Motion)
কোন বস্তুর ত্বরণ একটি নির্দিষ্ট বিন্দু থেকে এর সরণের সমানুপাতিক এবং সর্বদা ঐ বিন্দু অভিমুখী হলে তাকে সরল দোলন গতি বলে।

সমীকরণে ঋণাত্মক চিহ্নটি নির্দেশ করে যে বল ও সরণ সর্বদা পরস্পর বিপরীত দিকে ক্রিয়া করে, অর্থাৎ বস্তু সাম্যাবস্থান থেকে যেদিকেই সরুক, বল তাকে সাম্যাবস্থানের দিকে ফিরিয়ে আনে। |F| বা |a| এর মান সরণের মানের সমানুপাতিক, কিন্তু দিক সর্বদা বিপরীত। এ বল F একটি প্রত্যায়নী বল।

যে বল কোনো সাম্যাবস্থান থেকে সরণকৃত বস্তুকে সর্বদা সাম্যাবস্থানের দিকে ফিরিয়ে আনে তাকে প্রত্যায়নী বল বলে। এই বল সর্বদা সরণের বিপরীত দিকে ক্রিয়া করে। যখন প্রত্যায়নী বল সরণের সমানুপাতিক হয়, তখন সরল দোলন গতি সৃষ্টি হয়।
